% Le résumé et l'abstract.
%
% TROIS RÈGLES GOUVERNENT CES DEUX PAGES.
%
% Elles n'introduisent AUCUN résultat neuf : tout ce qu'elles citent est établi dans un chapitre,
% et s'y trouve avec sa source.
%
% Elles n'écrivent AUCUN nombre en clair. Les valeurs passent par le registre des chiffres, comme
% partout ailleurs dans le rapport. Ces appels ne sont pas comptés par le contrôle du registre, qui
% ne lit que les chapitres et les annexes : tous les identifiants employés ici sont donc déjà
% employés ailleurs, et aucun ne dépend de cette page pour ne pas devenir orphelin.
%
% Elles disent en une phrase que les données sont générées. Un résumé qui l'omettrait laisserait
% un lecteur pressé croire à des mesures réelles.
%
% L'abstract est une traduction fidèle et non une variante : mêmes appels de registre, mêmes
% mots-clés. Son vocabulaire technique suit l'usage anglophone du domaine.

\chapter*{Résumé}
\addcontentsline{toc}{chapter}{Résumé}

Le service d'accueil et d'admission est le point d'entrée administratif d'un hôpital. Il ouvre les
dossiers, enregistre les rendez-vous, admet et oriente les patients, facture et recouvre. Chacun de
ces gestes laisse une trace dans le système d'information hospitalier. Le règlement intérieur confie
à ce service l'établissement des statistiques, et la nomenclature nationale fixe ce qu'il doit
remonter. Entre les données saisies et les indicateurs attendus, il manquait un chemin.

Aucune donnée réelle ne sort de l'établissement, car la confidentialité des données de santé
l'interdit. \textbf{La chaîne repose donc sur un jeu de données généré par un programme écrit pour
ce projet.} Ce jeu est calibré sur des volumes publiés par établissement dans le recueil statistique
national, et compte \chiffre{source-lignes-total} lignes sur \chiffre{source-tables} tables, sur
\chiffre{periode-jours} jours.

La chaîne va du fichier déposé jusqu'à l'écran. Elle compte \chiffre{modeles-total} modèles de
transformation, met en quarantaine les lignes refusées et garde l'historique des fiches patient. Un
modèle probabiliste rapproche les identités, et il est évalué contre une vérité terrain connue. Un
tableau de bord de \chiffre{tdb-pages} pages restitue \chiffre{tdb-indicateurs} indicateurs, chacun
avec sa définition et la décision qu'il sert.

Le résultat le plus directement utilisable est la correspondance avec la nomenclature. Sur
\chiffre{exigences-rubriques} rubriques attendues, \chiffre{exigences-calculables} sont calculables
et \chiffre{exigences-non-calculables} ne le sont pas, chacune faute d'un champ du système. Enfin,
\chiffre{exigences-sans-objet} sont sans objet pour l'établissement, et
\chiffre{exigences-classement-mixte} reçoit un classement mixte.

\textbf{Aucune grandeur produite par la chaîne n'est une mesure de l'établissement.} Le travail
montre que la chaîne sait produire ces grandeurs, et dire pour chacune si elle est calculable, ce
qui l'en empêche et ce qu'elle vaut.

\section*{Mots-clés}

Système d'information hospitalier ; chaîne de données ; données synthétiques ; entrepôt
dimensionnel ; rapprochement d'identités ; qualité des données ; tableau de bord.

\begin{otherlanguage}{english}
\chapter*{Abstract}
\addcontentsline{toc}{chapter}{Abstract}

The admissions department is a hospital's administrative front door. It opens patient records,
books appointments, admits and directs patients, issues invoices and collects payment. Every one of
these actions leaves a trace in the hospital information system. Hospital regulations give this
department the duty of producing statistics, and the national nomenclature sets out what it must
report. Between the data recorded and the indicators expected, a path was missing.

No real data leaves the hospital, because the confidentiality of health data forbids it.
\textbf{The pipeline therefore rests on synthetic data produced by a program written for this
project.} The data are calibrated on volumes published per hospital in the national statistical
yearbook, and hold \chiffre{source-lignes-total} rows across \chiffre{source-tables} tables, over
\chiffre{periode-jours} days.

The pipeline runs from the landed file to the screen. It comprises \chiffre{modeles-total}
transformation models, quarantines rejected rows and keeps the history of patient records. A
probabilistic model links patient identities, and it is evaluated against known ground truth. A
\chiffre{tdb-pages}-page dashboard presents \chiffre{tdb-indicateurs} indicators, each with its
definition and the decision it serves.

The most directly usable result is the mapping to the nomenclature. Of
\chiffre{exigences-rubriques} required items, \chiffre{exigences-calculables} are computable and
\chiffre{exigences-non-calculables} are not, each for lack of a field in the system. Finally,
\chiffre{exigences-sans-objet} do not apply to this hospital, and
\chiffre{exigences-classement-mixte} receives a mixed classification.

\textbf{No figure this pipeline produces measures the hospital's actual activity.} The work shows
that the pipeline can produce these figures, and state for each whether it is computable, what
prevents it, and what it is worth.

\section*{Keywords}

Hospital information system; data pipeline; synthetic data; dimensional warehouse; record linkage;
data quality; dashboard.
\end{otherlanguage}
